\documentclass[10pt,letterpaper]{article}
\usepackage{cornell-notes}

\title{Annex C.05.09: Clause 19 Diagnostics Library}
\author{}
\date{}
\setDocTitle{Annex C.05.09: Clause 19 Diagnostics Library}
\setDocSubtitle{C++ 2024}
\setDocOwner{C++ 2024 Cornell Notes}
\setDocDate{\today}
\setCornellCollection{C++ 2024 Cornell Notes}
\setCornellUnitType{Annex}
\setCornellUnitNumber{C.05.09}
\setCornellUnitTitle{Clause 19 Diagnostics Library}
\hypersetup{pdftitle={Annex C.05.09: Clause 19 Diagnostics Library},pdfsubject={Cornell notes},pdfauthor={}}

\begin{document}
\maketitle
\begin{CornellOverviewBox}{Overview}
\textbf{Classification:} Informative annex\\[2pt]
\textbf{Learning objective:} Understand the role and technical guidance associated with Clause 19: diagnostics library.
\end{CornellOverviewBox}\begin{CornellSummaryBox}{Summary}
{\sffamily\bfseries\color{Accent} Summary}\par\vspace{3pt}Section C.5.9 focuses on Clause 19: diagnostics library. The section supplies the rules and semantic distinctions needed to interpret this topic in context with its cross-references. This material is informative: it supports understanding and navigation but does not independently create a conformance obligation.
\end{CornellSummaryBox}

\end{document}
